% !TeX root = ../main.tex
\section{Related Games}
\label{sec:related}

Tactical Formation is a folk game, not a game designed on the basis of the games
below. We compare them in order to situate its mechanics, without claiming
any historical relationship between the games. Folk games in different
places often share mechanics with one another.

\subsection{Games with a jump move}

In Checkers~\citep{schaeffer2007checkers}, a piece jumps over an
\emph{opponent's} piece to capture it, and the jumped-over piece is removed
from the board. In Halma and its relative Chinese Checkers, a piece may jump
over \emph{any} adjacent piece, its own or an opponent's, and the jumped-over
piece is \emph{not} removed~\citep{britannica_halma}. Chinese Checkers is
played by several players and has served as a test domain for multi-player
search algorithms~\citep{nijssen2013mcts}. The jump in Tactical Formation matches Halma
on both of these points. The differences are that Tactical Formation allows
\emph{only} jumps (there is no single step into an adjacent square), a turn
contains a single jump, and the jump triggers the sandwich-flip rule. Halma
has no capture at all.

\subsection{Games that capture by conversion}

In Reversi/Othello, a player \emph{places a new piece} and flips every
opponent piece sandwiched in a straight line, in \emph{all} directions at
once. The Logistello program of \citet{buro1997othello} once defeated the
world champion, and Othello $8\times8$ has recently been proven a
\emph{draw} under optimal play by both sides~\citep{takizawa2023othello}.
Rolit is a Reversi variant for up to 4 players, and it has been used as a
test domain for Paranoid, max$^n$ and Best-Reply
Search~\citep{schadd2011brs,nijssen2013mcts}. The
Reversi-like class of sandwich-flip games therefore already has a
multiplayer version. The sandwich rule of Tactical Formation is close to Reversi, with
three differences: pieces are \emph{moved} rather than placed; when several
sandwich lines are available, the mover takes only \emph{one}; and the total
number of pieces never changes.

The idea that ``capture is a change of side and the piece count is
conserved'' is not new either. In \emph{Shoulder-Pole Chess} from Quang Nam, a piece
steps along the lines of a 25-point board. If it moves into the point
between two enemy pieces in a straight line, those two pieces change side
(capture by \emph{intervention}: the moving piece ``carries'' the pair like a
shoulder pole, which gives the game its name), and a piece surrounded until it has no move also changes
side~\citep{mytour_coganh,fptshop_coganh}. The capture geometry of Shoulder-Pole Chess
is the reverse of Tactical Formation: in Shoulder-Pole Chess one's own piece moves \emph{between}
two enemy pieces, whereas in Tactical Formation enemy pieces are caught \emph{between}
two of one's own. In \emph{Ming Mang} (Tibet), pieces slide in straight lines
like a rook, and a line of enemy pieces sandwiched in a straight line passes
to the capturing side~\citep{wiki_mingmang}. This is the closest precedent
for the sandwich-conversion of Tactical Formation. \emph{Ataxx} also has a jump move,
after which every enemy piece \emph{adjacent} to the landing square changes
side, not along a line; its one-step move creates a new piece, so the total
piece count grows~\citep{gebop_ataxx}.

\subsection{Other piece-moving games}

\emph{Lines of Action}, designed by Claude Soucie and described in the
collection of \citet{sackson1969gamut}, also moves existing pieces instead
of placing new ones; the winning objective is to connect one's pieces, and
a captured piece is still removed from the board, as in Checkers. \emph{Focus} (also called \emph{Domination}) by
Sid Sackson~\citep{sackson1969gamut} keeps captured pieces so that they can
return to the board, so it too treats capture as transformation rather than
removal.

\subsection{Recording and digitising traditional games}

The Ludii system~\citep{piette2020ludii} describes board games, including
many traditional games, in a common game description language, so that
machines can play and compare them. This paper follows the same spirit: it
records a traditional game, models it, and implements it in order to
preserve it. Adding Tactical Formation to Ludii is a natural next step
(Section~\ref{sec:conclusion}).

\subsection{Overview comparison}

Table~\ref{tab:comparison} summarises the similarities and differences. To
our knowledge, no recorded game has all five of the following features of
Tactical Formation at once: (1)~movement only by jumping over any piece, with the
jumped-over piece staying in place; (2)~sandwich capture along a straight
line that converts pieces, but in a single chosen direction; (3)~a
two-layer vertex--centre board, with 8 directions for vertex pieces and 4
for centre pieces; (4)~a constant total number of pieces; (5)~play for 2 to
4 players. Each feature on its own already appears in other games.

\begin{table}[H]
  \centering
  \caption{Comparison of Tactical Formation with related games.}
  \label{tab:comparison}
  \footnotesize
  \renewcommand{\arraystretch}{1.2}
  \begin{tabularx}{\linewidth}{l*{6}{>{\centering\arraybackslash}X}}
    \toprule
    \textbf{Feature} & \textbf{Checkers} & \textbf{Halma / Chinese Checkers}
      & \textbf{Reversi / Rolit} & \textbf{Shoulder-Pole Chess} & \textbf{Ming Mang}
      & \textbf{Tactical Formation} \\
    \midrule
    Movement         & Diagonal step, jump & Step or jump & Place new piece
                     & Single step & Rook-like slide & Jump only \\
    Jumps over       & Enemy piece & Any piece & --- & --- & --- & Any piece \\
    Jumped-over piece & Removed & Stays & --- & --- & --- & Stays \\
    Capture          & Jump & None & Sandwich, all directions
                     & Inter\-vention, surround & Sandwich & Sandwich, one direction \\
    Captured piece   & Removed & --- & Changes side & Changes side
                     & Changes side & Changes side \\
    Total pieces     & Decreases & Constant & Increases & Constant & Constant
                     & Constant \\
    Players          & 2 & 2--6 & 2 / 2--4 & 2 & 2 & 2--4 \\
    Position layers  & 1 & 1 & 1 & 1 & 1 & 2 (vertex + centre) \\
    \bottomrule
  \end{tabularx}
\end{table}
